\section{Introduction}
Next-scale visual representations decompose an image into complete coarse token maps followed by progressive residual refinements at finer scales~\cite{var}. Because each scale spans the full latent field and conditions the next finer scale, it provides both a progressively refinable description and a generative prior for untransmitted content.

Under a fixed per-image wireless symbol budget, however, transmitting only complete scale maps leaves frame capacity unused whenever the next complete residual map does not fit---even though a substantial prefix of its tokens would.

Prior progressive codecs such as ARPC truncate only at complete scale boundaries~\cite{arpc}. Ada-TokenCom jointly adapts 1D token prefix length and the modulation and coding scheme (MCS) under a long-term average symbol budget~\cite{adatokencom}, while ResARC encodes extra residuals to compensate for suffix-generation errors~\cite{resarc}. Deep joint source--channel coding (JSCC) methods such as SwinJSCC map images directly to channel symbols and serve as system-level baselines~\cite{swinjscc}. What remains open is how to terminate within a residual scale under a strict per-frame symbol budget, while keeping every candidate prefix independently decodable and exploiting the recovered boundary tokens at the receiver.

This letter proposes NeST-Com. Its main contribution pairs an actual-length frame-fitting rule---selecting the longest independently terminated within-scale token prefix that fits the protected frame body via a coder-specific nondecreasing-length property---with a prefix-constrained receiver that folds recovered boundary tokens into VAR's known-token overwrite for both residual correction and finer-scale conditioning. Experiments on 500 ImageNet images isolate the gains from capacity utilization and boundary-token overwrite, and quantify the resulting quality--computation tradeoff.
